% First-failure event tree; q_t is conditional on survival.
\begin{tikzpicture}[x=1mm,y=1mm,font=\figurefont,
  state/.style={draw=BookBlue,fill=BookPaper,rounded corners=1pt,
    minimum width=20mm,minimum height=9mm,align=center,inner sep=1mm},
  failure/.style={state,draw=BookAmber},
  flow/.style={-{Stealth},thick,draw=BookInk}]
  \node[state] (root) at (12,0) {存活 $S_0$};
  \node[failure] (e0) at (54,15) {首次失败\\$E_0$};
  \node[state] (s1) at (54,-12) {存活 $S_1$};
  \node[failure] (e1) at (98,0) {首次失败\\$S_1\cap E_1$};
  \node[state] (s2) at (98,-26) {存活 $S_2$};
  \draw[flow,BookAmber] (root) -- node[above] {$q_0$} (e0);
  \draw[flow] (root) -- node[below] {$1-q_0$} (s1);
  \draw[flow,BookAmber] (s1) -- node[above] {$q_1$} (e1);
  \draw[flow] (s1) -- node[below] {$1-q_1$} (s2);
  \node[align=left,anchor=west,text width=64mm] at (1,-32)
    {两步整回合事件：\\$F_2=E_0\cup E_1$\\概率 $1-(1-q_0)(1-q_1)$};
  \draw[BookRule] (1,-45) -- (110,-45);
  \node[anchor=west] at (1,-52) {训练期另对多个回合取并集：};
  \node[state] (run1) at (16,-65) {回合1\\$F_H^{(1)}$};
  \node[state] (run2) at (54,-65) {回合2\\$F_H^{(2)}$};
  \node[state] (runk) at (96,-65) {回合$K$\\$F_H^{(K)}$};
  \draw[flow] (run1) -- (run2);
  \draw[flow,dashed] (run2) -- node[above] {$\cdots$} (runk);
  \node[anchor=west,text=BookMuted] at (1,-78)
    {$F_{\text{训}}=\bigcup_{k=1}^K F_H^{(k)}$，策略可能随训练改变};
\end{tikzpicture}
